\chapter{Replication}
\label{ch:replication}

\section{Replication}\label{replication}

We primarily consider \emph{remote} replication and assume a SAN-based
environment.

\subsection{Terminology}\label{terminology}

\begin{description}
\item[Production host]
in normal operations runs applications that access \emph{source} data.
\item[Source]
or production LUNs (volumes) hold data normally accessed by the
production host.
\item[Target]
LUN(s) where production data is replicated
\item[Recoverability]
is the restoration of data from replicas to source.
\item[Restartability]
is the restart of business operations using the replicas.
\end{description}

\subsection{Point-in-time vs continuous
replica}\label{point-in-time-vs-continuous-replica}

\begin{description}
\item[Point-in-time]
is a snapshot of the source LUN at some particular time.
\item[Continuous]
replica is in sync with the production data under normal operation.
\end{description}

\subsection{Replica uses}\label{replica-uses}

\begin{description}
\item[Alternative backup source]
from a PIT replica to avoid impacting on production performance.
\item[Fast recovery]
on source device failure: restoring data to different new source device
or restarting production directly on the replica.
\item[Reporting or analytics]
undertaken on replica to avoid impacting production performance.
\item[Testing platform]
for testing changes / updates on a testing environment before applying
to production environment.
\end{description}

\section{Modes}\label{modes}

There are two replication modes, synchronous and asynchronous.

\subsection{Synchronous}\label{synchronous}

shows the general sequence of synchronous replication.

\subsection{Asynchronoous}\label{asynchronoous}

\section{Technologies}\label{technologies}

\subsection{Host-based}\label{host-based}

\subsection{Array-based}\label{array-based}

\subsubsection{Synchronous}\label{synchronous-1}

\subsubsection{Asynchronous}\label{asynchronous}

\subsection{Network-based}\label{network-based}
